\section{Related Work}

\paragraph{Benchmarking symbolic regression.}
The symbolic regression community has repeatedly emphasized that scientific progress is constrained by weak benchmarking discipline as much as by algorithmic limitations. The original SRBench effort of \citet{lacava2021srbench} established a large, reproducible benchmark for comparing symbolic regression methods across both real-world and ground-truth problems under a common evaluation pipeline. More recently, the benchmark has been extended and revisited as a living community resource, with explicit discussion of hyperparameter tuning, execution constraints, energy-aware reporting, and benchmark maintenance policies \citep{aldeia2025call}. Our work is aligned with this line of thinking, but addresses a different layer of the problem: rather than only expanding the benchmark, we build an operational substrate that helps heterogeneous methods and datasets become executable under a shared protocol in the first place.

\paragraph{Datasets for scientific symbolic regression.}
Dataset design has become a central issue in symbolic regression research. The SRSD benchmark of \citet{matsubara2022rethinking} argues that many existing symbolic regression datasets are poorly matched to the scientific-discovery setting and proposes more realistic formula sampling ranges together with tree-aware equation comparison. More recently, \citet{shojaee2025llmsrbench} introduced LLM-SRBench to evaluate scientific equation discovery under settings that reduce trivial memorization for language-model-based systems. These works highlight the importance of split design, target semantics, and formula-aware evaluation. SR-Workbench builds on this perspective by standardizing multiple benchmark families into a common dataset schema with explicit validation, ID-test, OOD-test, metadata, and optional formula validation.

\paragraph{Unified or hybrid symbolic regression frameworks.}
Several works attempt to unify symbolic regression capabilities at the algorithmic level. A representative example is the unified deep symbolic regression framework of \citet{landajuela2022udsr}, which combines multiple symbolic regression solution strategies within a single modular model family and shows strong performance on SRBench. Such efforts are important because they demonstrate that hybridization across search strategies can be effective. However, their primary objective is algorithmic unification inside a model framework. Our objective is different: we unify independently maintained repositories, environments, datasets, and execution semantics across methods that were not originally designed to coexist in one experimental stack.

\paragraph{Automated research infrastructures.}
Recent automated research systems have started to address the broader problem of turning papers, repositories, and experiments into partially automated optimization loops. AutoSOTA \citep{li2026autosota}, for example, formulates an end-to-end research automation pipeline spanning paper grounding, environment repair, experiment execution, ideation, and validity supervision. Our work is inspired by this infrastructure perspective, but is deliberately lighter and more domain-specific. Instead of aiming at a general multi-agent scientific automation system, SR-Workbench focuses on symbolic regression as a concrete research vertical with its own requirements: equation outputs, complexity-sensitive evaluation, OOD reporting, and optional ground-truth formula validation. In this sense, our contribution is a symbolic-regression-specific substrate on which more ambitious automated loops can later be built.

\paragraph{Positioning of this work.}
Taken together, prior work shows that symbolic regression needs both better benchmarks and better method design, while automated research systems suggest a path toward reducing repetitive experimental overhead. SR-Workbench occupies the space between these directions. It is not a new symbolic regression model, and it is not yet a fully autonomous research agent. Instead, it provides a unified, auditable operating layer that connects heterogeneous methods, heterogeneous datasets, and bounded iterative improvement under one symbolic-regression-specific workflow.
